\section*{Indice degli Argomenti: Teoria dei Segnali (Modulo: Segnali Determinati e Aleatori)}

\begin{itemize}

    \item \textbf{Introduzione e Classificazione dei Segnali}
          \begin{itemize}
              \item Definizione di segnale e importanza nelle telecomunicazioni
              \item Classificazione: segnali determinati e aleatori
              \item Dominio e codominio: segnali a tempo continuo/discreto, ampiezza continua/discreta (segnali analogici, quantizzati, numerici)
              \item Conversione analogico/digitale
          \end{itemize}

    \item \textbf{Analisi dei Segnali nel Dominio del Tempo}
          \begin{itemize}
              \item Energia e Potenza media di un segnale
              \item Classificazione in segnali di energia e segnali di potenza
              \item Valor medio e Valore efficace (RMS)
          \end{itemize}

    \item \textbf{Trasformate di Fourier Generalizzate}
          \begin{itemize}
              \item Segnale gradino unitario
              \item Impulso unitario di Dirac ($\delta$) e sue proprietà (proprietà campionatrice, ritardo, dualità)
              \item Trasformata di Fourier di segnali periodici (spettro a righe)
              \item Prima formula di somma di Poisson
          \end{itemize}

    \item \textbf{Analisi dei Segnali Periodici (Serie di Fourier)}
          \begin{itemize}
              \item Segnali sinusoidali, fasori e formule di Eulero
              \item Sviluppo in Serie di Fourier: forma polare reale, forma complessa e forma rettangolare
              \item Spettro di ampiezza e di fase (unilatero e bilatero)
              \item Proprietà della Serie di Fourier (linearità, simmetria hermitiana, segnali pari e dispari)
              \item Ricostruzione approssimata del segnale e Fenomeno di Gibbs (ripple)
              \item Esempi notevoli: onda quadra, onda triangolare, treno di impulsi rettangolari (segnale rect), raddrizzatore a doppia semionda
          \end{itemize}

    \item \textbf{Campionamento dei Segnali}
          \begin{itemize}
              \item Teorema del campionamento (Nyquist-Shannon)
              \item Frequenza e intervallo di campionamento
              \item Fenomeno del ripiegamento dello spettro (Aliasing)
          \end{itemize}

    \item \textbf{Sistemi Monodimensionali a Tempo Continuo e Filtri}
          \begin{itemize}
              \item Definizione di sistema e proprietà fondamentali: linearità, stazionarietà, causalità, memoria, stabilità (BIBO), invertibilità
              \item Sistemi Lineari Stazionari (SLS) e Risposta impulsiva $h(t)$
              \item Integrale di Convoluzione e sue proprietà
              \item Risposta in frequenza $H(f)$
              \item Sistemi in cascata e in parallelo
              \item Filtri ideali: passa-basso, passa-alto, passa-banda, elimina-banda
              \item Filtri reali, circuiti RC/LC e Criterio di Paley-Wiener
              \item Distorsione lineare di ampiezza e di fase introdotta dai filtri
              \item Misure logaritmiche: il Decibel (dB) e banda a -3dB
          \end{itemize}

    \item \textbf{Analisi Spettrale: Energia, Potenza e Correlazione}
          \begin{itemize}
              \item Densità spettrale di energia e Teorema di Parseval
              \item Densità spettrale di potenza
              \item Funzione di autocorrelazione e autocovarianza
              \item Teorema di Wiener-Khintchine (relazione tra autocorrelazione e spettro)
              \item Banda al 99\% dell'energia
          \end{itemize}

    \item \textbf{Processi Aleatori e Rumore}
          \begin{itemize}
              \item Definizione di processo aleatorio, spazio campione e funzioni campione
              \item Processi parametrici e variabili aleatorie
              \item Caratterizzazione statistica di I e II ordine (valor medio statistico, varianza, autocorrelazione e autocovarianza)
              \item Processi stazionari in senso stretto e in senso lato (WSS)
              \item Processi ergodici (uguaglianza tra media temporale e media statistica)
              \item Processi aleatori Gaussiani
              \item Filtraggio di un segnale aleatorio
              \item Modelli di rumore: Rumore bianco e Rumore termico (di Johnson)
          \end{itemize}

    \item \textbf{Modulazioni Analogiche in Banda Passante}
          \begin{itemize}
              \item Scopo della modulazione e definizioni (modulante, portante, modulato)
              \item Modulazioni di Ampiezza (AM):
                    \begin{itemize}
                        \item DSB-AM (Banda laterale doppia con portante) e indice di modulazione
                        \item DSB-SC (Banda laterale doppia a portante soppressa)
                        \item SSB (Banda laterale unica: USSB e LSSB) tramite filtraggio o Trasformata di Hilbert
                        \item VSB (Banda laterale vestigiale)
                    \end{itemize}
              \item Tecniche di demodulazione: Rivelatore di inviluppo vs Demodulazione coerente/sincrona
              \item Modulazioni d'Angolo:
                    \begin{itemize}
                        \item Modulazione di Fase (PM)
                        \item Modulazione di Frequenza (FM)
                        \item Deviazione di frequenza/fase e Regola di Carson per il calcolo della banda
                    \end{itemize}
              \item Multiplazione a divisione di frequenza (FDM) e bande di guardia
          \end{itemize}

\end{itemize}